\documentclass[11pt]{article}

\usepackage[margin=1in]{geometry}
\usepackage{amsmath}
\usepackage{booktabs}
\usepackage[hidelinks]{hyperref}
\usepackage{microtype}

\title{Bulk-Silicon Wannier Construction and\\Optimization-Diagnosis Research Protocol}
\author{Project Koios}
\date{Calculation-owned protocol}

\begin{document}
\maketitle

\begin{abstract}
This protocol separates scientific target-space design, local numerical optimization,
arithmetic or toolchain instability, and downstream operator validation in the
bulk-silicon Wannier route.  It prevents failed localization from being mislabeled as
either proof of physical infeasibility or evidence of floating-point instability, and
it prevents successful Wannier90 convergence from being promoted directly to scientific
acceptance.  Exact thresholds and evidence dispositions remain calculation-owned.
\end{abstract}

\section{Scope and authority}

The protocol applies after an authenticated plane-wave parent, ordered NSCF Bloch frame,
Wannier90 neighbor topology, and AMN/MMN/EIG interface have been established.  Reusable
software owns parsing, identity checks, bounded execution, matrix invariants, and lineage.
The calculation owns rank, projections, windows, convergence coordinates, observable
tolerances, candidate selection, and scientific acceptance.

A trial projection initializes a Wannier gauge.  It does not impose a Slater--Koster
Hamiltonian, hopping shell, or fitted parameter set.  Slater--Koster reduction begins
only after acceptance of the full Wannier operator.

\section{Distinct failure classes}

The following classes shall not be conflated.

\subsection{Target-space or construction failure}

Rank, parent-band inventory, outer and frozen windows, projections, reciprocal mesh, and
neighbor topology define the represented state space and initialization.  Sensitivity to
these coordinates is scientific-construction evidence.  A construction can create a
poorly conditioned or strongly nonconvex optimization landscape without implying
arithmetic failure.

\subsection{Local-optimizer failure}

Wannier90 uses a projection-initialized local optimization; it does not provide a
certificate of the global minimum of the spread functional.  Iteration exhaustion,
large late spread excursions, or materially different gauges from different optimizer
coordinates establish nonstationarity or local-optimizer sensitivity.  They do not prove
that no acceptable gauge exists.

\subsection{Arithmetic or toolchain instability}

Arithmetic or toolchain instability requires direct evidence such as non-finite values,
corrupt or inconsistent artifacts, failed matrix identities beyond declared tolerances,
same-input nondeterminism, or controlled sensitivity to arithmetic library or executable
choice.  Optimizer excursions alone are insufficient for this diagnosis.

\subsection{Scientific inadequacy}

A converged and structurally valid Wannier occurrence can still be scientifically
inadequate.  Interpolation, valley identity and location, edge energy, derivatives,
curvature or effective mass, and real-space locality are separate acceptance channels.

\section{Mandatory diagnostic procedure}

For every localization occurrence:

\begin{enumerate}
  \item Require exact executable and input identities, a bounded authorization record,
    a fresh or explicitly declared restart state, and no silent retry.
  \item Assess disentanglement and Wannierisation separately using their native
    convergence messages and criteria.  A zero return code is not a substitute.
  \item Verify finite dimensions and k-point ordering, gauge unitarity,
    disentanglement isometry when applicable, Hamiltonian representative pairing, and
    Hermiticity.
  \item Inspect convergence trajectories.  Distinguish smooth contraction toward a
    declared threshold from intermittent plateaus, renewed excursions, and budget
    exhaustion.
  \item Hold authenticated interface artifacts fixed when diagnosing optimizer or
    target-space changes.  Record every changed scientific and algorithmic coordinate.
  \item Permit only a prospectively bounded optimizer study.  If the trajectory remains
    irregular and observably gauge-sensitive, revise the scientific construction rather
    than performing an open-ended optimizer search.
  \item After native convergence, perform construction- and withheld-mesh observable
    tests before accepting the operator.
\end{enumerate}

\section{Global-minimum policy}

Scientific acceptance does not require an unsupported claim that a global spread
minimum was found.  A locally converged candidate may be accepted only if it is
structurally valid, reproducible under the declared protocol, and satisfies all
observable and locality tolerances.  Competing-minimum studies become necessary when two
or more converged initializations or constructions produce materially different target
observables or locality conclusions.  Comparisons shall use gauge-invariant spectra and
within-candidate locality metrics unless an explicit alignment map exists.

The lowest final spread among different target spaces is not, by itself, a selection
rule because changing the frozen manifold changes the represented optimization problem.

\section{Prospective convergence policy}

For prospective $C_{6\text{-}B20\text{-}W21}$ localization, the calculation adopts
the convergence-derived tuple from the precise protocol in pinned
\texttt{aiida-wannier90-workflows} 3.1.0.  Its builder multiplies the per-atom tolerance
meta-parameters by the number of sites.  For the exact two-atom primitive cell this
gives
\begin{align}
  \texttt{conv\_tol} &= 2.0\times10^{-9}\ \text{\AA}^{2}, &
  \texttt{conv\_window} &= 6, \\
  \texttt{dis\_conv\_tol} &= 2.0\times10^{-9}, &
  \texttt{num\_cg\_steps} &= 200,
\end{align}
with \texttt{num\_iter}=40000 and \texttt{dis\_num\_iter}=40000.  The emitted
\texttt{dis\_conv\_tol} remains dimensionless; ``per atom'' describes the workflow's
scaling convention.  Both native convergence messages remain mandatory, and the
40000-step values are fixed eligibility ceilings rather than evidence of a global
minimum or scientific adequacy.

This adaptation does not adopt the precise protocol's factor-three band count,
automatic 0.1~\AA$^{-1}$ mesh rule, PBEsol pseudopotential family, projections, energy
windows, MPI policy, or restart handlers.  Historical occurrences retain the exact
policies under which they ran and are not reclassified retrospectively.  No automatic
extension beyond the adopted ceilings is permitted.

\section{Current silicon evidence}

The all-conduction-frozen rank-10 construction froze eight conduction states and no
valence states at every source k point.  Its localization trajectories remained
nonstationary through as many as 4000 iterations, including 277 post-1000 spread changes
larger than $0.05\ \text{\AA}^{2}$ and a final large excursion at iteration 3934.  A
500/1000 diagnostic found close source-mesh conduction spectra but material dense-grid,
valley, and curvature differences.  These observations support optimization-landscape
and target-space sensitivity, not arithmetic instability.

The successor retained the executable, rank, projections, mesh, and byte-identical
AMN, MMN, EIG, and NNKP interface artifacts, but changed the frozen interval to
$[-5.85,8.00]$~eV in the native QE energy reference.  This freezes all four valence
states and zero to two sampled low-energy conduction states per k point.  With
\texttt{dis\_num\_iter = 2000}, disentanglement converged at iteration 1290 and
Wannierisation converged at iteration 189.  The final total spread was
$19.532839107\ \text{\AA}^{2}$; retained gauge-unitarity and disentanglement-isometry
residuals were approximately $1.58\times10^{-10}$ and $1.32\times10^{-10}$, and the
maximum Hamiltonian Hermiticity residual was $0.0$~eV.

Because the same authenticated interface and executable admit convergence after the
target-space revision, the supported diagnosis is that the earlier obstruction arose
primarily from construction-induced optimization geometry.  This is not a proof that
the new occurrence is a global spread minimum or a scientifically accepted operator.

The subsequent no-calculator characterization verified source reconstruction and
recorded the candidate's interpolation, valley, curvature, and locality behavior without
assigning pass/fail status to candidate-only observations.  A prospectively frozen
408-point direct-parent NSCF then completed the intended-use comparison.  All six
candidate-to-parent valley displacements, $0.02120$--$0.03086$~\AA$^{-1}$, met the
$0.034465$~\AA$^{-1}$ target.  The candidate indirect gap was $0.52261$~eV versus the
direct-parent $0.61125$~eV, a 14.50\% relative error against the 5\% target.  On the
frozen candidate-centered primary stencils, all transverse-mass comparisons and two of
six longitudinal-mass comparisons exceeded the 8\% target.  The candidate therefore
remains provider-eligible and structurally valid but is not adequate for the declared
SSSP-pilot intended use.

\section{Construction-mesh successor recommendation}

The negative intended-use result ends optimizer tuning for this candidate.  It shall not
be repaired by retrospective tolerance changes, additional localization iterations, or
mixing optimizer changes with a scientific-construction revision.  Component-operator
construction and Slater--Koster reduction remain blocked.

The first successor shall isolate reciprocal construction-mesh density.  Candidate
$C_6$ uses a Gamma-origin $6\times6\times6$ ordered mesh while holding fixed the accepted
$S_1$ structure and parent, QE 42/336~Ry cutoffs, fixed occupations,
\texttt{nosym} and \texttt{noinv} set true, \texttt{conv\_thr} set to
$1.0\times10^{-12}$~Ry, full-accuracy Davidson diagonalization, 12 parent bands,
rank 10, non-spinorial character, outer window
$[-5.85,19.02]$~eV, frozen window $[-5.85,8.00]$~eV, and the two projection lines
\texttt{Si:sp3} and \texttt{Si:l=0:r=2}.  The localization controls also remain
\texttt{dis\_num\_iter} 2000, \texttt{dis\_conv\_tol} $1.0\times10^{-10}$,
\texttt{dis\_conv\_window} 3, \texttt{dis\_mix\_ratio} 0.2,
\texttt{num\_iter} 4000, \texttt{conv\_tol} $2.0\times10^{-9}$~\AA$^2$,
\texttt{conv\_window} 6, enabled \texttt{precond}, \texttt{num\_cg\_steps} 200,
native-default \texttt{trial\_step} 2.0, unset \texttt{fixed\_step}, and disabled
\texttt{guiding\_centres}.

The $C_6$ route consists of distinct NSCF, Wannier90 preprocessing,
\texttt{pw2wannier90.x}, and localization occurrences.  Each calculator invocation
requires its own exact preflight and bounded authorization.  No occurrence in this
protocol is automatically authorized.

The first exact $C_6$ localization occurrence retained
\texttt{dis\_num\_iter}=2000 and completed natively, but exhausted that budget at an
absolute fractional change of $4.679\times10^{-8}$ against the unchanged dimensionless
$1.0\times10^{-10}$ criterion.  Wannierisation subsequently converged at iteration 316,
but cannot repair the unconverged subspace.  This is an optimizer-budget failure and
produces no provider-eligible operator.  The monotonically decreasing changes over
iterations 1800--2000 support preparation of one fresh bounded eligibility test with
\texttt{dis\_num\_iter}=4000.  Every target-space coordinate and all other optimizer
controls remain fixed, and no failed checkpoint or output may be consumed.  This budget
extension is an explicit optimizer-coordinate amendment, not evidence that the original
occurrence converged or that either local optimization reaches a global minimum.  The
fresh occurrence subsequently converged Disentanglement at iteration 3068, but
Wannierisation did not meet its six-consecutive-step spread-change gate by iteration
4000.  Its final absolute spread change was $1.26\times10^{-5}$~\AA$^2$ against the
unchanged $2.0\times10^{-9}$~\AA$^2$ criterion, and the late trajectory remained
nonmonotonic with large excursions.  The provider therefore rejected the output.  The
prospective stop rule now applies: no further automatic budget extension follows and
$C_8$ shall not be selected merely because $C_6$ lacks an eligible operator.

\subsection{Target-space candidate \texorpdfstring{$C_{6\text{-}T1}$}{C6-T1}}

The authenticated $C_6$ parent eigenvalues span
$[-5.8341319732,20.8036050385]$~eV.  The inherited outer window
$[-5.85,19.02]$~eV contains 12 parent states at 142 construction points, 11 at 62,
and only the rank-minimum 10 at 12; 86 states across 74 points lie above its maximum.
The fixed numerical window therefore does not define a mesh-stable parent-band
inventory.  This observation does not prove the cause of either optimizer failure, but
it ends pure mesh-only attribution and motivates a separate target-space candidate.

Candidate $C_{6\text{-}T1}$ changes only the outer window to
$[-6.0,21.0]$~eV.  On the authenticated 216-point interface this defines the outer
subspace as the span of all 12 computed parent states at every construction point, with
finite lower and upper margins of 0.1658680268 and 0.1963949615~eV.  The frozen window,
rank, projections, mesh, source order, energy reference, interface bytes, and all
optimizer controls remain fixed, including \texttt{dis\_num\_iter}=4000 and
\texttt{num\_iter}=4000.  The occurrence shall start fresh and consume no checkpoint or
failed localization output.  Because the window does not change NNKP topology or the
QE-generated EIG, AMN, and MMN data, those authenticated interface artifacts may be
reused without new preprocessing or conversion.

$C_{6\text{-}T1}$ is not a continuation of the one-variable mesh test.  Its prepared
localization request remains unexecuted and is not selected as the next occurrence after
the parent-band review below.

\subsection{AiiDA-moderate parent-band candidate
\texorpdfstring{$C_{6\text{-}B20}$}{C6-B20}}

The next selected target-space stage adopts the band-count heuristic, but not the
complete protocol, from \texttt{aiida-wannier90-workflows} release 3.1.0 at pinned
commit \texttt{1c0115cb7c77}.  The pinned moderate
workflow sets \texttt{num\_bands\_factor}=2.0, and its estimator takes the maximum of
the factor-scaled occupied-state and projection counts and the corresponding counts
plus four.\footnote{\href{https://github.com/aiidateam/aiida-wannier90-workflows/blob/1c0115cb7c77541a0bd549bc05005df435d651b6/src/aiida_wannier90_workflows/workflows/protocols/base/wannier90.yaml\#L6-L25}{Pinned moderate-protocol source};
\href{https://github.com/aiidateam/aiida-wannier90-workflows/blob/1c0115cb7c77541a0bd549bc05005df435d651b6/src/aiida_wannier90_workflows/utils/pseudo/__init__.py\#L292-L342}{pinned band-estimator implementation}.}
For the non-spin-polarized two-atom cell, eight pseudopotential valence electrons, and
ten declared trial projections, this gives
\begin{equation}
  N_{\mathrm{bands}}=\max(8,8,20,14)=20.
\end{equation}
For this parent-band decision, the calculation-owned adaptation cites the AiiDA
workflow only for that factor-two inventory rule.  The separately declared precise
convergence adaptation controls prospective localization and does not alter the band
count.  Neither adaptation adopts the workflow's default PBEsol pseudopotential,
derives projections through its pseudo machinery, or claims that AiiDA executed or
validated this campaign.

Candidate $C_{6\text{-}B20}$ changes only the NSCF parent-band count from 12 to 20.
The accepted structure and parent, PBE pseudopotential, 42/336~Ry cutoffs, ordered
Gamma-origin $6\times6\times6$ mesh, occupations, $10^{-12}$~Ry electronic threshold,
full-accuracy Davidson solver, and symmetry controls remain fixed.  Rank 10 remains the
prospective represented dimension; 20 calculated bands do not imply rank 20.

One separately authorized $C_{6\text{-}B20}$ NSCF completed successfully.  Its QEXSD
contains 216 points and 20 finite bands spanning $[-5.8341319732,30.2720123939]$~eV,
and its saved state verifies.  The ordered points and weights match the earlier
12-band $C_6$ occurrence exactly, while its first 12 eigenvalues agree within
$4.14\times10^{-11}$~eV.  Under the inherited $[-5.85,19.02]$~eV window, the enlarged
inventory contains 10--14 states per point with histogram
$\{10\!:\!12,11\!:\!62,12\!:\!96,13\!:\!24,14\!:\!22\}$.  Under the prepared
$[-6.0,21.0]$~eV window it contains 12--14 states with histogram
$\{12\!:\!48,13\!:\!108,14\!:\!60\}$.  Bands 13 and 14 therefore enter both
investigated windows; band 15 remains above 21~eV everywhere.

These observations authenticate the broader parent inventory.  The subsequent
calculation-owned target-space decision selects candidate
$C_{6\text{-}B20\text{-}W21}$ with outer window $[-6.0,21.0]$~eV.  This places the
upper boundary 0.1964~eV above the observed band-12 maximum and 0.2572~eV below the
observed band-15 minimum, while admitting 12--14 states and therefore at least two
outer-space dimensions beyond rank 10 at every point.  The frozen window, rank, and
projection family remain unchanged.  This selection is an optimizer-eligibility
target-space redesign, not evidence of convergence, a global minimum, or scientific
adequacy.

Bands 13--20 require new interface artifacts; the existing 12-band NNKP, EIG, AMN, and
MMN data cannot be reused.  Fresh preprocessing, conversion, and localization remain
distinct occurrences with separate authorization.  The AiiDA precise factor of 3.0
would give 30 bands under the same adaptation, but $C_{6\text{-}B30}$ is a conditional
upper diagnostic and does not follow automatically from failure of
$C_{6\text{-}B20\text{-}W21}$.

The revealed 408-point parent set may be used for successor development and residual
attribution, but it is no longer withheld for $C_6$ or any later candidate.  The
existing 5\% indirect-gap, 3\% Gamma--X valley-position, and 8\% per-mass targets remain
unchanged development criteria.  If $C_6$ meets every development criterion, its exact
operator and predictions shall be frozen before any new parent values are observed.  If
$C_6$ does not meet them, no $8\times8\times8$ successor follows automatically.  The
operator shall review the channel-resolved $C_4$-to-$C_6$ changes: a separately declared
$C_8$ mesh study is appropriate only when the errors contract consistently enough to
support a mesh-density hypothesis.  Otherwise rank, band inventory, windows, and
projections return to separate scientific review and shall not be changed together.

A $12\times12\times12$ construction, denoted $C_{12}$, is a bounded upper mesh
diagnostic rather than an automatic successor.  The $C_4$, $C_6$, $C_8$, and $C_{12}$
meshes contain 64, 216, 512, and 1728 points, respectively, so $C_{12}$ is 27 times the
original source-point inventory.  It becomes scientifically eligible for a separate
preflight only when a completed $C_8$ review still supports the mesh-density hypothesis
and either the intended-use criteria remain unmet or material mesh sensitivity remains.
No direct jump from $C_4$ to $C_{12}$ is part of this protocol.

Passing the development and prospective confirmation criteria at $C_{12}$ would
establish adequacy only for that finite represented operator and declared intended use.
A stronger claim of construction-mesh convergence requires either comparison with a
finer declared mesh or a prospective real-space tail/aliasing bound showing that omitted
Fourier content cannot materially change the target observables.  Mesh size alone does
not supply that bound, and $C_{12}$ cannot repair an inadequate band inventory, rank,
window, or projection family.

A future confirmation set shall remain disjoint from the $C_4$, $C_6$, possible $C_8$,
and possible $C_{12}$ construction meshes and from the revealed 408-point set.  Its geometry is fixed
prospectively as follows.  Six symmetry-projected centers lie on the signed Cartesian
axes at the mean revealed parent-valley radius
$0.9629174643168844$~\AA$^{-1}$, equal to $0.8381689694043609$ of Gamma--X.  Around each
center, use the right-handed rotated Cartesian basis
\begin{equation}
 U=\begin{pmatrix}
 1/\sqrt{3} &  1/\sqrt{2} &  1/\sqrt{6} \\
 1/\sqrt{3} & -1/\sqrt{2} &  1/\sqrt{6} \\
 1/\sqrt{3} &  0           & -2/\sqrt{6}
 \end{pmatrix}.
\end{equation}
For each center, include the center, all six axial and twelve mixed-sign quadratic
points at each of $0.002$, $0.006$, and $0.012$~\AA$^{-1}$, and the twelve axial capture
points at $0.024$ and $0.048$~\AA$^{-1}$.  After periodic de-duplication the intended
inventory contains 402 points.  Exact fractional coordinates, construction-disjointness,
and point roles must be generated and authenticated before use.

This confirmation geometry is parent-informed because its centers use the already
revealed parent minima; it shall not be described as blind to geometry.  Its direct
parent eigenvalues nevertheless remain unobserved.  A successor candidate's 402-point
predictions, fitting conventions, and unchanged intended-use criteria must be frozen
before a separately authorized direct-parent occurrence reveals those values.

\section{Successor acceptance sequence}

For $C_6$ and any separately approved successor:

\begin{enumerate}
  \item verify exact source-mesh spectral reconstruction and parent-eigenvalue
    correlation;
  \item evaluate dense interpolation, valley identity and location, edge energies,
    derivatives, curvature/effective masses, and real-space locality on retained
    development data;
  \item freeze the candidate, candidate predictions, fitting conventions, and criteria
    before observing confirmation parent values;
  \item evaluate the 402-point confirmation set without changing the frozen criteria;
    and
  \item only after full-operator acceptance, begin component-operator construction and
    a separately declared Slater--Koster reduction.
\end{enumerate}

Failure in any channel shall be attributed to the applicable parent-model,
discretization, subspace, gauge, interpolation, or truncation error class.  Passing
software checks or native convergence alone shall not be reported as scientific
validation.

\section{Nonclaims}

This protocol does not certify a global minimum, physical accuracy, pseudopotential
transferability, mesh convergence, effective masses, locality, or a Slater--Koster
model.  It authorizes no calculator invocation, evidence publication, release, or
retrospective acceptance.

\end{document}
